%% kapa: exam=1 round=26 year=2015 period=2 form=A subject=appraisal_law subject_ko=부동산관계법규 seq=22 number=22
%% source: Q-Net 감정평가사 기출문제 (cst00309) artlSeq=5138001 file=extracted/1008/5138001/제26회_감정평가사 1차_기출/제26회_감정평가사 1차_기출/2교시_감평_A형.pdf
%% conversion: kordoc 4.24.51 -> tools/gen1.py (auto); review=proofread-sonnet; flags=none
\begin{question}{22}
부동산가격공시 및 감정평가에 관한 법령상의 내용으로 옳지 않은 것은?
\begin{choices}
  \item 법원에 계속중인 소송 또는 경매를 위한 토지등의 감정평가는 감정평가업자의 업무에 속한다.
  \item 금고 이상의 형의 선고유예를 받고 그 선고유예기간 중에 있는 자는 감정평가사가 될 수 없다.
  \item 감정평가업자는 정당한 사유없이 그 업무상 알게 된 비밀을 누설하여서는 아니된다.
  \item 감정평가업자는 감정평가서의 원본을 교부일부터 3년 이상, 감정평가서의 관련 서류를 교부일부터 1년 이상 보존하여야 한다.
  \item 감정평가협회는 법인으로 한다.
\end{choices}
\end{question}
\answer{4}
